\mode*
\section{Composition: pipes, redirection, scripts}
\label{sec:composition}

The third educational objective is what makes the shell more than a clumsy
file manager: composing programs with pipes and redirection, and freezing a
composition into a script.

\subsection{Documented difficulties}

Composite command production is exactly where
\textcite{Doane1992UnixPrompts} located the difficulty: users could know
the component commands yet fail to compose them, and even targeted help
prompts did not suffice for novices.
% TODO(#1): The systematic phase should search for later replications or
% CS-education studies of pipeline construction difficulties.

\subsection{Candidate critical aspects}

% XXX Preliminary; to be consolidated once the review phase is complete.
\begin{itemize}
  \item \emph{Uniform interfaces}: programs read text in and write text
    out, which is why any producer can feed any consumer.
  \item \emph{The pipe is not the file}: redirection to a file versus
    piping to a program.
  \item \emph{A script is nothing new}: the same commands, replayable
    (cf.\ the history-to-script demonstration used in the course).
\end{itemize}

\subsection{Patterns of variation}

\begin{itemize}
  \item Contrast (sink varies, producer invariant): the same command's
    output to the screen, to a file (\texttt{>}), to a program
    (\texttt{|}).
  \item Generalisation (pipeline shape invariant, programs vary):
    \texttt{X | sort}, \texttt{X | wc -l} across different \texttt{X}.
  \item Fusion: build a pipeline solving a real task, then turn the
    history into a script and rerun it.
\end{itemize}
